\section{Conditional probability}

The previous section explained the idea before the formula: when information is
received, the relevant space of reasoning is restricted. If we know that an
event \(B\) occurred, then outcomes outside \(B\) are no longer compatible with
what we know. If we then ask about another event \(A\), the favourable outcomes
are those in
\[
A\cap B.
\]

We now turn this idea into a precise definition. The definition should not be
read as an arbitrary rule. It is the mathematical expression of two decisions:
first restrict to the information event, then rescale so that the information
event becomes the new whole space.

\subsection*{The finite uniform case first}

Let \(\Omega\) be a finite sample space in which all elementary outcomes are
equally likely. Suppose we are told that an event \(B\) occurred, and assume
\(B\) is not empty. After this information, the possible outcomes are the
elements of \(B\).

If we ask for the probability of \(A\) after knowing \(B\), then only the
outcomes in \(B\) matter. Among these, the favourable outcomes are precisely the
elements of
\[
A\cap B.
\]
Thus, in the finite uniform case, the natural value is
\[
\frac{|A\cap B|}{|B|}.
\]

This expression is not guessed. It is the ordinary uniform probability formula,
but applied in the restricted space \(B\) instead of the original space
\(\Omega\).

\begin{examplebox}
A fair die is rolled. Let
\[
A=\{4,5,6\}
\]
be the event that the result is greater than \(3\), and let
\[
B=\{2,4,6\}
\]
be the event that the result is even. If we know \(B\), then the relevant space
is \(\{2,4,6\}\). Inside this space, the event \(A\) occurs for \(4\) and
\(6\). Therefore the probability of \(A\) after knowing \(B\) is
\[
\frac{2}{3}.
\]
This is the ratio \(|A\cap B|/|B|\), because
\[
A\cap B=\{4,6\}.
\]
\end{examplebox}

\subsection*{From counting to probability mass}

The expression
\[
\frac{|A\cap B|}{|B|}
\]
works only when the restricted elementary outcomes are equally likely. In a
general probability model, outcomes may have different probabilities. Then
counting outcomes is not enough. We must measure probability mass.

The information event \(B\) has probability \(P(B)\). This is the total mass of
all outcomes compatible with the information. The event \(A\cap B\) has
probability \(P(A\cap B)\). This is the part of that mass for which the event
\(A\) also occurs.

After learning \(B\), the whole region \(B\) becomes the new reference space. It
must therefore have total probability \(1\) in the updated situation. The part
of \(B\) where \(A\) occurs should occupy the same relative share that it had
inside \(B\) before rescaling. That relative share is
\[
\frac{P(A\cap B)}{P(B)}.
\]

This gives the definition.

\begin{definitionbox}
Let \(A\) and \(B\) be events in a probability space, and assume
\[
P(B)>0.
\]
The \textbf{conditional probability of \(A\) given \(B\)} is
\[
P(A\mid B)=\frac{P(A\cap B)}{P(B)}.
\]
\end{definitionbox}

The condition \(P(B)>0\) is essential. If the information event has probability
zero, then it cannot be made into a new probability-one reference region by
ordinary rescaling.

\subsection*{Why the denominator is \(P(B)\)}

The denominator is not a technical detail. It represents the whole amount of
probability still compatible with the information. Once we know \(B\), the event
\(B\) must become certain in the new calculation:
\[
P(B\mid B)=1.
\]
The definition gives exactly this:
\[
P(B\mid B)=\frac{P(B\cap B)}{P(B)}=rac{P(B)}{P(B)}=1.
\]

Similarly, an event impossible under the information should have conditional
probability zero. If \(A\cap B=\varnothing\), then
\[
P(A\mid B)=\frac{P(\varnothing)}{P(B)}=0.
\]
This matches the idea that, once \(B\) is known, an event disjoint from \(B\)
cannot occur.

\subsection*{Conditional probability as a new probability assignment}

For a fixed event \(B\) with \(P(B)>0\), the rule
\[
A\mapsto P(A\mid B)
\]
defines a new probability assignment on the original events, but interpreted
under the information \(B\).

It has the usual properties. First,
\[
P(A\mid B)\geq 0,
\]
because probabilities are non-negative. Second,
\[
P(\Omega\mid B)=\frac{P(\Omega\cap B)}{P(B)}=rac{P(B)}{P(B)}=1.
\]
Third, if \(A_1\) and \(A_2\) are disjoint, then
\[
P((A_1\cup A_2)\mid B)
=
\frac{P((A_1\cup A_2)\cap B)}{P(B)}.
\]
Since
\[
(A_1\cup A_2)\cap B=(A_1\cap B)\cup(A_2\cap B),
\]
and the two pieces on the right are disjoint, additivity gives
\[
P((A_1\cup A_2)\mid B)
=
\frac{P(A_1\cap B)+P(A_2\cap B)}{P(B)}.
\]
Therefore
\[
P((A_1\cup A_2)\mid B)=P(A_1\mid B)+P(A_2\mid B).
\]

This confirms that conditional probability is not a separate kind of arithmetic.
It is ordinary probability after changing the reference space.

\subsection*{Multiplication rule}

The definition can be rewritten as
\[
P(A\cap B)=P(A\mid B)P(B).
\]
This is called the multiplication rule. It says: to get the probability that
both \(A\) and \(B\) occur, first measure the probability of the information
region \(B\), and then multiply by the probability that \(A\) occurs inside that
region.

The same intersection can also be read in the opposite order:
\[
P(A\cap B)=P(B\mid A)P(A),
\]
provided \(P(A)>0\).

Thus
\[
P(A\mid B)P(B)=P(B\mid A)P(A).
\]
This identity will later become the basis of Bayes' theorem.

\begin{remarkbox}
The multiplication rule is not an additional assumption. It is the definition of
conditional probability rearranged. Its usefulness comes from the fact that many
multi-stage situations are naturally described by first entering one event and
then asking for another event inside it.
\end{remarkbox}

\subsection*{Example: drawing two cards without replacement}

Draw two cards from a standard deck without replacement, and record the ordered
sequence. Let
\[
A=\{\text{the first card is an ace}\},
\]
and
\[
B=\{\text{the second card is an ace}\}.
\]
We want to understand the probability that both cards are aces.

First,
\[
P(A)=\frac4{52}.
\]
If \(A\) has occurred, then one ace has already been removed. There are now
\(51\) cards left, of which \(3\) are aces. Therefore
\[
P(B\mid A)=\frac3{51}.
\]
By the multiplication rule,
\[
P(A\cap B)=P(A)P(B\mid A)=\frac4{52}\cdot\frac3{51}.
\]

This calculation is not merely a formula. It follows the physical structure of
the experiment: first card, then second card after the first card has changed the
deck.

\subsection*{Example: medical-style testing without Bayes yet}

Suppose a population is divided into two groups: people who have a certain
condition and people who do not. Let
\[
C=\{\text{the person has the condition}\},
\]
and let
\[
T=\{\text{the test result is positive}\}.
\]
The conditional probability
\[
P(T\mid C)
\]
means: among people with the condition, what proportion test positive?

The conditional probability
\[
P(C\mid T)
\]
means: among people with a positive test result, what proportion have the
condition?

These are different questions. They use different reference spaces. The first
uses \(C\) as the information event. The second uses \(T\) as the information
event. There is no reason for the two numbers to be equal.

This distinction is one of the main reasons conditional probability must be
introduced carefully. The vertical bar in \(P(A\mid B)\) is not punctuation. It
separates the event being measured from the information that defines the new
space of reasoning.

\subsection*{Conditional probabilities in tables}

Conditional probability can often be read from a table. Suppose a group of
students is classified by whether they passed an exam and whether they attended
review sessions:
\[
\begin{array}{c|cc|c}
 & \text{passed} & \text{did not pass} & \text{total}\\
\hline
\text{attended} & 36 & 9 & 45\\
\text{did not attend} & 24 & 31 & 55\\
\hline
\text{total} & 60 & 40 & 100
\end{array}
\]
Let
\[
A=\{\text{student passed}\},
\]
and
\[
B=\{\text{student attended}\}.
\]
Then \(P(A\mid B)\) is not computed from all \(100\) students. The information
\(B\) restricts attention to the \(45\) students who attended. Among them,
\(36\) passed. Thus
\[
P(A\mid B)=\frac{36}{45}.
\]

By contrast,
\[
P(B\mid A)=\frac{36}{60},
\]
because now the information is that the student passed, and among the \(60\)
students who passed, \(36\) attended. Again, the two conditional probabilities
answer different questions.

\subsection*{What this definition teaches}

The definition
\[
P(A\mid B)=\frac{P(A\cap B)}{P(B)}
\]
should be remembered through its meaning:

\begin{itemize}
    \item \(B\) is the information that becomes the new reference space;
    \item \(A\cap B\) is the part of that reference space where the event of
    interest also occurs;
    \item division by \(P(B)\) rescales the information region so that it has
    total probability \(1\).
\end{itemize}

\begin{theorembox}
Conditional probability is ordinary probability after the space of reasoning has
been restricted by information. The formula records this restriction and the
necessary rescaling.
\end{theorembox}
